\documentclass[11pt]{article}
\usepackage[a4paper,margin=25mm]{geometry}
\usepackage{amsmath,amssymb,amsthm}
\newtheorem{theorem}{Theorem}
\title{Structural certificates for the $\mathrm{Sp}(4,5)$ atlas table}
\author{Private L054 finite application}\date{}
\begin{document}\maketitle
\begin{theorem}
Relative to the acquired numerical catalogue and structural registry,
all 307 rows of the current companion table headed $\mathrm{Sp}(4,5)$ have
the displayed structural descriptions. The ambient group is
$\operatorname{PSp}_4(5)$, of order 4680000, as constructed in the acquired
numerical execution. Each extension description asserts an actual
normal subgroup and quotient, with splitting or nonsplitting as displayed;
it is not a unique abstract classification by an extension label.
\end{theorem}

\paragraph{Source and actual groups.}
The table is attributed to Thomas Connor and Dimitri Leemans and is dated
7 June 2013. All fourteen native pages were directly inspected. The source's
$\mathrm{Sp}(4,5)$ notation is used for the projective group in this table;
the group of all form-preserving matrices has twice the displayed order.
The precise named input is the original completed GAP construction
$G:=\operatorname{PSp}(4,5)$. Its saved ordered ambient generators agree
literally with the accepted robust numerical catalogue. Every proper
subgroup is saved by literal 156-point generators, with no constructor or
other printed group name substituted for an embedding.

The numerical source bijection and weighted inclusions are acquired, and
each structural occurrence is checked against its exact source row,
order and embedded generators. The source punctuation is retained. In
particular row 10 is $2\mathbin{\cdot}(A_5\times A_5)$ and row 30 is
$Q_8\mathbin{\cdot}(A_5\times3)$. Parentheses around the single atoms
$Q_{20},D_{20},D_{10}$ in rows 18,19,31,33,133 are redundant; their original
text is preserved. The $5^3$ and $5^{1+2}$ atoms are distinct. The earlier
dicyclic conventions for $Q_{12},Q_{20}$ and unsigned extraspecial
conventions are acquired unchanged.

\paragraph{Finite proof interface.}
There are 196 acquired concrete seeds and 55 new models: 38 split, five
direct, eleven nonsplit, and one elementary abelian model. Every factor
reference is acyclic. All finite groups are closed from their literal
generators. An identity-preserving bijection which respects right
multiplication by each ordered generator is a homomorphism by induction
on words, hence an isomorphism. This acquired generator-edge criterion
verifies each factor and occurrence map; equal orders alone never supply
an identification.

For a split model, literal checks establish $N\triangleleft H$, a subgroup
$K$, $N\cap K=1$, and $|N||K|=|H|$. Thus $H=NK$. A colon includes an
actual noncommuting generator pair; a direct product has elementwise
commuting factors. The new elementary abelian model has commuting
prime-order basis generators and the exact required order. Row 1 acquires
the named ambient construction and its full-order subgroup identity.

\newpage
\paragraph{Nonsplitting certificates.}
The literal reader uses only previously accepted exhaustive criteria.
For a prime-order quotient, it checks every outside element: a complement
would supply an element of that prime order. For a quotient $S_3$, six
verified cosets give quotient generators $xN,yN$ and the relations
$x^2=y^3=(xy)^2=1$. Every pair of lifts is tested; a pair satisfying these
relations would generate a complement, and a complement supplies such a
pair. Both directions follow from the presentation and quotient order.

For the central-kernel model, the reader verifies $N\leq Z(H)$ and
$x,y\in H$ with $1\ne[x,y]\in N$. Their quotient images commute, whereas
all pairs of lifts have the same nonidentity commutator. A complement is
therefore impossible. The actual quotient is separately identified by its
literal coset action and a certified registry map. This is the previously
reviewed commuting-lifts criterion, retaining its original contribution.

Two models use the acquired general Cayley-lift criterion, which permits
a nonabelian kernel or quotient. Let $Q\cong H/N$ be an identified concrete
model with ordered generators $q_1,\ldots,q_r$, and let $x_i$ be actual
lifts. Enumerate every tuple in
\[
 (Nx_1)\times\cdots\times(Nx_r).
\]
For each tuple, follow the complete right-generator graph of $Q$, assigning
the identity first and multiplying by the chosen lifts. A consistent
assignment on every edge gives a homomorphism $s:Q\to H$. Its composition
with the quotient map is the identified isomorphism, so $s$ is injective
and its image is a complement. Conversely any complement supplies one
of these tuples and a consistent assignment. Thus one failed edge for
every tuple proves nonsplitting. The reader verifies complete disjoint
cosets, the ordered generator lifts, the quotient identification, and
all $|N|^r$ choices. A failed software complement search is used only to
propose a certificate, never as its proof.

\paragraph{Validation and boundaries.}
The first bounded proposal pass stopped at an omitted candidate quotient
case after 28 positive rows. Its failed receipt and completed witnesses
are retained. A successor changed only that candidate case and visited
the 278 unfinished rows; all returned positive with clean error output.
The source trees and accepted mathematical criteria were unchanged.

Independent literal replay verifies all model and occurrence certificates
in ordinary and optimized modes. The reader also checks GAP's four physical
numeric-array continuations at bracket boundaries without permitting
continuations inside JSON strings. Focused damages test source assignments,
actual embeddings, punctuation, hashes, quotient lifts and parsing. The
acquired Cayley criterion's small complete-subgroup controls are included.

No new ambient census, unique interpretation of extension notation,
historical PDF-byte identity, whole-source admission, question resolution,
formalization or measurement is asserted. The larger source-associated
tables remain in scope until individually discharged.
\end{document}
